\documentclass[11pt,a4paper]{article}

\usepackage[utf8]{inputenc}
\usepackage[T1]{fontenc}
\usepackage{amsmath,amssymb,amsthm}
\usepackage{geometry}
\usepackage{booktabs}
\usepackage{xcolor}
\usepackage{physics}
\usepackage{hyperref}
\usepackage{enumitem}
\usepackage{tikz}
\usetikzlibrary{arrows,positioning}

\geometry{margin=1in}

\newtheorem{theorem}{Theorem}[section]
\newtheorem{proposition}[theorem]{Proposition}
\theoremstyle{definition}
\newtheorem{definition}[theorem]{Definition}

\newcommand{\Tc}{\mathcal{T}}
\newcommand{\dgeo}{\delta}
\newcommand{\Mpl}{M_{\text{Pl}}}
\newcommand{\murad}{\mu_0}
\newcommand{\nurad}{\nu_0}
\newcommand{\mH}{m_H}
\newcommand{\Esyn}{E_{\text{syn}}}

\title{\textbf{QSD Supplement S2: Energy Scale Derivation}\\[0.5em]
\large From Planck Mass to Electroweak Scale}

\author{Marcus Hubka\\
\small Independent Researcher, Czech Republic}

\date{February 2026}

\begin{document}

\maketitle

\begin{abstract}
This supplement derives the QSD energy scale hierarchy, showing how the electroweak scale emerges from the Planck mass through information-theoretic reduction. We derive the factor $2^{-60}$ from the code state space, the geometric enhancement by $\varphi^2$, and the strain correction $(1-6\dgeo)$.
\end{abstract}

\tableofcontents

\section{The Hierarchy Problem}

The Standard Model has two scales: $\Mpl = 1.22 \times 10^{19}$ GeV and $v_{\text{EW}} \approx 246$ GeV, a ratio of $\sim 5 \times 10^{16}$.

QSD resolves this through the master formula:
\begin{equation}
\boxed{\murad = \frac{\Mpl}{2^{60}} \times \varphi^2 \times (1 - 6\dgeo)}
\end{equation}

\section{The Information Reduction Factor $2^{-60}$}

\begin{theorem}[State Space Volume]
\begin{equation}
\mathcal{N} = 16^{15} = 2^{60} = 1.153 \times 10^{18}
\end{equation}
\end{theorem}

\begin{proof}
The $K_4$ Hilbert space has dimension $D = 2^4 = 16$. The number of independent interaction channels is $N = D^2 - 1 = 15$ (dimension of $\mathfrak{su}(4)$). Total: $D^N = 16^{15} = 2^{60}$.
\end{proof}

Base scale:
\begin{equation}
E_{\text{base}} = \frac{\Mpl}{2^{60}} = \frac{1.221 \times 10^{19}}{1.153 \times 10^{18}} = 10.59~\text{GeV}
\end{equation}

\section{Geometric Enhancement by $\varphi^2$}

The golden ratio emerges from $C_5$ geometry. The radion inherits:
\begin{equation}
\murad \propto E_{\text{base}} \times \varphi^2 = 10.59 \times 2.618 = 27.72~\text{GeV}
\end{equation}

\section{Strain Correction $(1 - 6\dgeo)$}

Each of 6 edges of $K_4$ loses fraction $\dgeo$ due to $K_4/C_5$ mismatch:
\begin{equation}
1 - 6\dgeo = 1 - 6(0.01432) = 0.914
\end{equation}

\section{Complete Calculation}

\begin{align}
\murad &= \frac{1.221 \times 10^{19}}{1.153 \times 10^{18}} \times 2.618 \times 0.914 \\
&= 10.59 \times 2.618 \times 0.914 = \boxed{25.34~\text{GeV}}
\end{align}

\section{Consistency Check}

Dimensionless VEV:
\begin{equation}
\nurad = \frac{\Tc}{\dgeo} = \frac{0.3607}{0.01432} = 25.18
\end{equation}

Agreement with $\murad$/GeV: 0.6\% --- nontrivial consistency.

\section{Derived Masses}

\begin{table}[h]
\centering
\begin{tabular}{lccc}
\toprule
\textbf{Quantity} & \textbf{Formula} & \textbf{QSD} & \textbf{Exp.} \\
\midrule
Radion & $\Mpl/2^{60} \cdot \varphi^2(1-6\dgeo)$ & 25.34 GeV & --- \\
Higgs & $5\murad$ & 126.7 GeV & 125.1 GeV \\
Syndrome & $\murad(1+8\dgeo)$ & 28.24 GeV & $\sim$28 GeV \\
\bottomrule
\end{tabular}
\end{table}

All from $\Mpl$ with zero free parameters.

\end{document}
